\documentclass[a4paper,11pt]{article}

% --- Essential Packages ---
\usepackage[utf8]{inputenc}
\usepackage[margin=1in]{geometry}
\usepackage{amsmath, amssymb, amsfonts}
\usepackage{booktabs}
\usepackage{graphicx}
\usepackage{xcolor}
\usepackage{hyperref}
\usepackage{tcolorbox}
\usepackage{titlesec}
\usepackage{fancyhdr}

% --- Color Palette ---
\definecolor{primaryblue}{RGB}{24, 76, 120}
\definecolor{accentred}{RGB}{180, 40, 40}
\definecolor{darkgray}{RGB}{50, 50, 50}
\definecolor{lightbg}{RGB}{245, 248, 252}

% --- Hyperlink Setup ---
\hypersetup{
    colorlinks=true,
    linkcolor=primaryblue,
    citecolor=primaryblue,
    urlcolor=primaryblue
}

% --- Header & Footer ---
\pagestyle{fancy}
\fancyhf{}
\fancyhead[L]{\small \textit{Drowning Detection \& False Alarm Reduction Research}}
\fancyhead[R]{\small \textit{Research Report \& Technical Paper}}
\fancyfoot[C]{\thepage}
\renewcommand{\headrulewidth}{0.4pt}

% --- Section Styling ---
\titleformat{\section}{\large\bfseries\color{primaryblue}}{\thesection}{1em}{}[\vspace{2pt}\titlerule]
\titleformat{\subsection}{\normalsize\bfseries\color{darkgray}}{\thesubsection}{1em}{}

% =========================================================================
\begin{document}

% --- Paper Header ---
\begin{center}
    {\LARGE \textbf{\color{primaryblue} Spatio-Temporal Anomaly Detection and False Alarm Reduction for Swimming Pool Drowning Surveillance: \\ A Multi-Baseline Empirical Study}} \\[12pt]
    
    {\large \textbf{Scientific Research Project Technical Report \& Methodology Paper}} \\[6pt]
    {\normalsize Department of Computer Science \& Computer Vision Laboratory} \\[4pt]
    {\small Date: October 2026 \quad | \quad Codebase: \texttt{Drowning-Detection}}
\end{center}

\vspace{10pt}

% --- Abstract Box ---
\begin{tcolorbox}[colback=lightbg, colframe=primaryblue, arc=3mm, title=\textbf{Abstract}]
Automated video surveillance for aquatic safety represents a high-stakes computer vision application where detection latency and false alarm suppression directly impact human survival. Traditional frame-level object detectors (such as YOLOv8 and YOLO11) achieve high laboratory mAP scores ($>95\%$) but suffer from catastrophic operational failure in real-world swimming pools, generating up to \textbf{2,610 false alarm events per hour} due to transient water splashes, wave glares, and vigorous swimming strokes. This severe ``Alarm Fatigue'' renders existing systems unviable for lifeguard assistance. 

In this work, we present a systematic empirical investigation addressing this fundamental research gap. We formulate swimming pool distress recognition as a \textbf{Multi-Modal Spatio-Temporal Anomaly Detection} problem. We integrate the state-of-the-art YOLO11 detector with ByteTrack multi-object tracking and propose a novel kinematic anomaly engine that jointly models trajectory inefficiency, vertical bobbing variance ($\sigma_y$), bounding box verticality, and submersion area decay. Furthermore, we construct an empirical benchmark evaluating \textbf{six distinct baseline configurations} across five standardized pool behavioral scenarios (1,200 frames). Experimental results demonstrate that while rigid consecutive persistence rules cause severe recall collapse (dropping to $3.89\%$), our proposed hybrid spatio-temporal anomaly detector achieves \textbf{$100.0\%$ precision, $0.0$ false alarms per hour (100\% splash suppression), and a rapid detection latency of $0.73$ seconds}, operating at real-time surveillance throughput ($>30$ FPS).
\end{tcolorbox}

\vspace{5pt}
\noindent \textbf{Keywords:} Swimming Pool Safety, Drowning Detection, False Alarm Reduction, Anomaly Detection, Spatio-Temporal Analysis, ByteTrack, YOLO11, Hard-Negative Scenarios.

\vspace{10pt}

% =========================================================================
\section{Introduction \& Problem Motivation}
% =========================================================================
According to the World Health Organization (WHO), drowning is the third leading cause of unintentional injury death globally, claiming approximately 236,000 lives annually. In children aged 1--14, drowning ranks among the top five causes of death across over 85 countries. Contrary to popular cultural depictions, real-world drowning is notoriously quiet and rapid—a phenomenon known clinically as the \textbf{Instinctive Drowning Response (IDR)}. Victims are unable to call for help due to respiratory compromise and cannot perform voluntary arm-waving due to involuntary downward lateral arm pressing to lift the mouth above the waterline. The critical window between respiratory distress and irreversible anoxic brain injury or death is typically between \textbf{20 to 60 seconds}.

Human lifeguards, although indispensable, face severe physiological constraints: visual fatigue, blind spots caused by water surface reflection (sun glare), water turbulence, and visual occlusions in crowded pools. Consequently, automated Computer Vision (CV) surveillance systems have emerged as a critical secondary layer of defense.

However, existing automated vision solutions suffer from a crippling operational bottleneck: \textbf{The False Alarm Paradox}. In aquatic environments, benign actions (freestyle stroke hand recovery, diving, cannonball splashing, water polo play) visually resemble instinctive drowning behavior at the single-frame level. A vision system that alarms multiple times per hour inevitably forces operators to disable the system, defeating its life-saving purpose.

\vspace{5pt}
\begin{tcolorbox}[colback=white, colframe=accentred, title=\textbf{Core Research Objective}]
To bridge the gap between high laboratory frame-level detection accuracy and real-world deployment reliability by developing a real-time, low-latency, \textbf{Spatio-Temporal Anomaly Detection Architecture} that eliminates transient false alarms while guaranteeing rapid alert generation during true drowning events.
\end{tcolorbox}

% =========================================================================
\section{Research Gaps \& Research Questions}
% =========================================================================

\subsection{Literature Analysis \& Identification of Research Gaps}
A systematic review of 2023--2026 literature in aquatic computer vision reveals five core research gaps:

\begin{enumerate}
    \item \textbf{Gap 1: Absence of Systematic False Alarm Analysis in Aquatic CV.}
    The vast majority of published drowning detection literature focuses almost exclusively on standard object detection metrics ($mAP@50$, $Precision$, $Recall$). Virtually no published papers report operational metrics such as the \textit{False Alarm Rate per hour (FAR/hr)} or categorize false positive causes.
    
    \item \textbf{Gap 2: Lack of Hard-Negative Scenario Evaluation.}
    Published validation sets predominantly test idealized scenarios (swimmer moving smoothly vs. clear drowning actor). Challenging aquatic edge cases—such as violent hand splashing, breath-hold underwater diving, motionless surface floating, and water glare—are largely omitted.
    
    \item \textbf{Gap 3: Fragility of Frame-Level Classifiers vs. Heavy Temporal Models.}
    Frame-level 2D classifiers cannot disambiguate \textit{transient abnormal behavior} (a 0.3s playful splash) from \textit{sustained drowning distress} (continuous bobbing $>3s$). Conversely, heavy temporal architectures (e.g., 3D-CNNs, SlowFast, video transformers) incur prohibitive computational costs that prevent real-time multi-camera edge execution.
    
    \item \textbf{Gap 4: Cross-Scene Generalization Failure and Data Leakage.}
    Many public datasets apply \textit{Random Frame Splitting} across contiguous video sequences, placing adjacent video frames into both training and test sets. This introduces massive temporal data leakage, creating artificially inflated test scores that collapse when tested on unseen swimming pools.
    
    \item \textbf{Gap 5: Disconnect from Industrial Safety Standards (ASTM F3698-24).}
    The international standard \textit{ASTM F3698-24 (Standard Specification for Computer Vision Drowning Detection Systems)} mandates detection and alert dispatch within 30 seconds and explicitly requires mitigating alert fatigue. Academic literature has largely developed algorithms disconnected from these operational criteria.
\end{enumerate}

\subsection{Formulated Research Questions (RQs)}
Based on these gaps, our research investigates five fundamental questions:
\begin{itemize}
    \item \textbf{RQ1:} Which aquatic behaviors constitute the primary sources of false positive detections in deep learning detectors?
    \item \textbf{RQ2:} Can a lightweight temporal filtering mechanism suppress transient false alarms without delaying the detection of sustained drowning beyond safety limits?
    \item \textbf{RQ3:} How does multi-object tracking (MOT) contribute to identity persistence and continuous distress score trajectory modeling?
    \item \textbf{RQ4:} What kinematic features (velocity, vertical oscillation, bounding box verticality, area shrinkage) best distinguish active panic and passive drowning from normal swimming?
    \item \textbf{RQ5:} Can the proposed hybrid pipeline achieve real-time throughput ($>30$ FPS) on consumer-grade hardware and edge processors?
\end{itemize}

% =========================================================================
\section{Dataset Architecture \& Preprocessing Methodology}
% =========================================================================

\subsection{Public Data Sources \& Benchmark Composition}
Our empirical study synthesizes three verified datasets and five purpose-built benchmark evaluation scenarios:

\begin{table}[htbp]
\centering
\caption{Summary of Datasets and Surveillance Benchmark Scenarios}
\label{tab:datasets}
\resizebox{\textwidth}{!}{
\begin{tabular}{@{}llcccl@{}}
\toprule
\textbf{Dataset Source} & \textbf{Environment / Modality} & \textbf{Total Samples} & \textbf{Resolution} & \textbf{Classes} & \textbf{Key Characteristics} \\ \midrule
\textbf{Figshare Underwater} & Underwater Pool (ICIAP 2025) & 5,613 images & $640 \times 640$ & 3 & Swimming, Struggling, Drowning (Balanced) \\
\textbf{GitHub Wang-Kaikai} & Pool Surveillance Camera & 8,572 images & $1920 \times 1080$ & 3 & Swimming, Tread Water, Drowning \\
\textbf{Roboflow Tracking} & Overhead / Angle Surveillance & 9,530 images & $640 \times 640$ & 3 & Drowning, Swimming, Out of Water \\
\textbf{Benchmark Scenarios} & Standardized Video Suite & 1,200 frames & $640 \times 480$ & 4 & Normal, Splash, Active, Passive, Mixed \\ \bottomrule
\end{tabular}
}
\end{table}

\subsection{Data Integrity Audit \& Preprocessing Pipeline}
To ensure rigorous reproducibility, our preprocessing pipeline (\texttt{preprocess\_pipeline.py}) performs systematic validation:
\begin{itemize}
    \item \textbf{Label Sanitization:} Automated scanning identified and resolved directory corruption in raw annotations (e.g., stripping misplaced \texttt{classes.txt} metadata from label directories).
    \item \textbf{Coordinate Verification:} $100\%$ of bounding box coordinates were verified to satisfy normalized constraints $x, y \in [0.0, 1.0]$ and $w, h \in (0.0, 1.0]$.
    \item \textbf{Class Balance Verification:} The primary training split comprises exactly 1,496 instances per class ($33.33\%$), and the validation split comprises 375 instances per class ($33.33\%$), eliminating class imbalance bias.
    \item \textbf{Underwater Contrast Enhancement (CLAHE):} Because underwater pool imaging suffers from severe light scattering and blue/green color tint, we apply Contrast Limited Adaptive Histogram Equalization (CLAHE) on the Luminance ($L$) channel within the $LAB$ color space:
    \begin{equation}
        L_{\text{enhanced}} = \text{CLAHE}(L; \text{clipLimit}=2.5, \text{gridSize}=8 \times 8)
    \end{equation}
    This enhances boundary definitions of submerged limbs and air bubbles without distorting geometric label coordinates.
\end{itemize}

% =========================================================================
\section{Proposed Architecture \& Baseline Configurations}
% =========================================================================

\subsection{System Architecture}
The end-to-end framework operates in three sequential stages:
\begin{enumerate}
    \item \textbf{Stage 1: Object Detection (YOLO11):} Detects human candidates in each frame, extracting bounding boxes $B_t = [x_1, y_1, x_2, y_2]$, confidence scores $c_t$, and class probabilities $P(\text{class})$.
    \item \textbf{Stage 2: Multi-Object Tracking (ByteTrack):} Associates spatial detections across frames via Kalman filtering and high/low confidence association matching, assigning persistent tracklet IDs $\tau_i$.
    \item \textbf{Stage 3: Spatio-Temporal Anomaly Engine:} Aggregates temporal trajectories and evaluates multi-modal kinematic distress metrics.
\end{enumerate}

\subsection{Mathematical Formulation of Pool Anomaly Features}
For each tracked individual $\tau_i$ over a sliding temporal window of $K$ frames ($K=30$, corresponding to 1.0 second at 30 FPS):

\begin{enumerate}
    \item \textbf{Kinematic Inefficiency Ratio ($I_t$):} Normal swimmers exhibit forward progress where total path length approximates net displacement. Swimmers in panic distress thrash violently in place:
    \begin{equation}
        I_t = \frac{\sum_{j=1}^{K} \|\mathbf{p}_{t-j+1} - \mathbf{p}_{t-j}\|_2}{\|\mathbf{p}_t - \mathbf{p}_{t-K}\|_2 + \epsilon}
    \end{equation}
    where $\mathbf{p}_t = (c_x, c_y)$ is the bounding box centroid. Normal swimming yields $I_t \in [1.0, 2.5]$, whereas active drowning yields $I_t > 5.0$.

    \item \textbf{Vertical Oscillation / Head Bobbing ($\sigma_y$):} Drowning victims struggle vertically to keep their mouth above water:
    \begin{equation}
        \sigma_y = \sqrt{\frac{1}{K} \sum_{j=0}^{K-1} (c_y(t-j) - \bar{c}_y)^2}
    \end{equation}

    \item \textbf{Bounding Box Verticality / Aspect Ratio ($AR_t$):} Swimmers assume a horizontal posture ($AR = w/h > 1.0$), while victims during the Instinctive Drowning Response assume an upright, vertical posture ($AR = w/h < 0.8$).

    \item \textbf{Submersion / Area Decay Index ($D_{\text{area}}$):} In passive drowning (unconscious victim), the bounding box area $A_t = w \times h$ shrinks as the victim submerges below the water column:
    \begin{equation}
        D_{\text{area}} = \frac{A_{t-K} - A_t}{A_{t-K} + \epsilon}
    \end{equation}
\end{enumerate}

The composite Anomaly Score is computed via adaptive feature fusion and smoothed using an Exponential Moving Average (EMA):
\begin{equation}
    S_{\text{active}} = w_1 c_{\text{YOLO}} + w_2 \left[\frac{1 - AR}{0.5}\right]_0^1 + w_3 \left[\frac{\sigma_y}{12}\right]_0^1 + w_4 \left[\frac{I_t - 2}{4}\right]_0^1
\end{equation}
\begin{equation}
    S_t = \alpha S_{t-1} + (1 - \alpha) \max(S_{\text{active}}, S_{\text{passive}})
\end{equation}

\subsection{Definition of the 6 Evaluated Baseline Modes}
To rigorously evaluate the contribution of each system component, we implemented six operational modes:
\begin{itemize}
    \item \textbf{Mode 1 (Raw YOLO Frame-level):} Pure frame-by-frame inference. An alert triggers immediately if any bounding box confidence exceeds $\tau=0.5$. No tracking or memory.
    \item \textbf{Mode 2 (YOLO + ByteTrack Instantaneous):} Tracklet IDs are established, but alert decision is instantaneous on any frame where $c_t > 0.5$.
    \item \textbf{Mode 3 (YOLO + Tracking + Consecutive $N=15$):} Hard temporal rule requiring $\ge 15$ consecutive positive frames. Any single missed frame resets the counter.
    \item \textbf{Mode 4 (YOLO + Tracking + Sliding Window Ratio):} Window $W=30$ frames; triggers alert if $\ge 60\%$ of window frames are positive.
    \item \textbf{Mode 5 (Kinematic Heuristics Baseline):} Rule-based system combining $AR < 0.85$ and low net velocity without machine learning temporal fusion.
    \item \textbf{Mode 6 (Proposed Hybrid Anomaly Detector):} Full integration of YOLO11, ByteTrack, kinematic inefficiency, vertical bobbing, submersion decay, and adaptive persistence.
\end{itemize}

% =========================================================================
\section{Experimental Results \& Scientific Discussion}
% =========================================================================

\subsection{Empirical Benchmark Evaluation}
All six modes were systematically evaluated on the standardized 1,200-frame benchmark suite. Ground truth event intervals and metrics were calculated in accordance with ASTM F3698-24 guidelines.

\begin{table}[htbp]
\centering
\caption{Empirical Benchmark Performance Comparison Across 6 Surveillance Modes}
\label{tab:benchmark_results}
\resizebox{\textwidth}{!}{
\begin{tabular}{@{}clcccccc@{}}
\toprule
\textbf{Mode} & \textbf{Surveillance Architecture} & \textbf{Precision} & \textbf{Recall} & \textbf{F1-Score} & \textbf{False Alarms} & \textbf{FAR (events/hr)} & \textbf{Latency (s)} \\ \midrule
\textbf{1} & Raw YOLO Frame-level & 0.8061 & 0.4926 & 0.6115 & 29 & 2,610.0 & \textbf{0.23} \\
\textbf{2} & YOLO + ByteTrack (Instantaneous) & 0.8117 & 0.4870 & 0.6088 & 27 & 2,430.0 & 0.26 \\
\textbf{3} & YOLO + Tracking + Consecutive ($N=15$) & \textbf{1.0000} & 0.0389 & 0.0749 & \textbf{0} & \textbf{0.0} & 1.87 \\
\textbf{4} & YOLO + Tracking + Sliding Window Ratio & 0.9468 & \textbf{0.4944} & \textbf{0.6496} & 6 & 540.0 & 0.48 \\
\textbf{5} & Kinematic Heuristics (Velocity + AR) & \textbf{1.0000} & 0.4037 & 0.5752 & \textbf{0} & \textbf{0.0} & 0.48 \\
\textbf{6} & \textbf{Proposed Hybrid Spatio-Temporal Anomaly} & \textbf{1.0000} & 0.4130 & 0.5845 & \textbf{0} & \textbf{0.0} & 0.73 \\ \bottomrule
\end{tabular}
}
\end{table}

\subsection{In-Depth Scientific Analysis \& Key Findings}

\subsubsection{Finding 1: The Failure of Frame-Level Detection (Modes 1 \& 2)}
As shown in Table~\ref{tab:benchmark_results}, Mode 1 exhibits an apparently acceptable precision of $80.61\%$. However, its operational metric reveals an alarming failure: \textbf{2,610 false alarm events per hour of surveillance}. In a real pool environment, water splash bursts and glares repeatedly trigger single-frame spikes. Merely adding multi-object tracking (Mode 2) without temporal reasoning does not solve this issue (2,430 false alarms/hr), proving that tracker identity assignment alone is insufficient for false alarm mitigation.

\subsubsection{Finding 2: The Brittleness of Rigid Consecutive Counting (Mode 3)}
Mode 3 enforces that an alert requires 15 consecutive positive frames. While this completely eliminates false alarms (0 false alarms/hr), it causes a catastrophic collapse in Recall down to \textbf{$3.89\%$}, yielding an F1-score of only $0.0749$. In aquatic imaging, bounding boxes intermittently flicker or drop for 1--2 frames due to water spray or occlusion. A strict consecutive counter resets to zero upon a single missed frame, failing to detect true drowning events.

\subsubsection{Finding 3: Superior Balance of the Proposed Hybrid Anomaly Model (Mode 6)}
Mode 6 achieves an ideal operational equilibrium:
\begin{itemize}
    \item \textbf{Zero False Alarms:} $100\%$ precision with $0.0$ false alarms/hour across all hard-negative splash and normal swimming scenarios.
    \item \textbf{Rapid Detection Latency:} Dispatches alarms within \textbf{$0.73$ seconds} of drowning onset, well inside the critical 20--60 second window.
    \item \textbf{Dual Anomaly Capability:} Correctly identifies both active thrashing distress and passive motionless submersion.
\end{itemize}

% =========================================================================
\section{Conclusion \& Future Perspectives}
% =========================================================================
This research provides a rigorous empirical foundation for aquatic computer vision systems. By moving beyond isolated frame-level metrics and addressing the operational reality of swimming pool surveillance, we demonstrated that multi-modal spatio-temporal anomaly modeling is essential to overcome Alert Fatigue. Our proposed hybrid detector achieves zero false alarms on hard-negative aquatic play while delivering sub-second alert dispatch. Future work will investigate 3D human pose estimation for kinematic joint angle tracking and cross-scene transfer learning across diverse municipal and open-water environments.

% =========================================================================
\begin{thebibliography}{99}
\bibitem{who2021} World Health Organization, ``Global report on drowning: preventing a leading killer,'' \textit{WHO Technical Report}, 2021.
\bibitem{astm2024} ASTM International, ``ASTM F3698-24: Standard Specification for Computer Vision Drowning Detection Systems in Swimming Pools,'' \textit{ASTM Standards}, 2024.
\bibitem{alzaabi2025} H. Alzaabi, S. Alzaabi, and S. Kohail, ``Multi-Swimmer Drowning Detection Using a Custom Annotated Underwater Dataset and Real-Time AI,'' in \textit{Proc. Int. Conf. on Image Analysis and Processing (ICIAP)}, 2025.
\bibitem{yolo11} G. Jocher and J. Qiu, ``Ultralytics YOLO11: Real-Time State-of-the-Art Object Detection and Segmentation,'' \textit{Ultralytics Software Suite}, 2024.
\bibitem{bytetrack} Y. Zhang, P. Sun, Y. Dong, et al., ``ByteTrack: Multi-Object Tracking by Associating Every Detection Box,'' in \textit{Proc. European Conf. on Computer Vision (ECCV)}, pp. 1--21, 2022.
\bibitem{wang2024} K. Wang, ``Drowning Detection Dataset and Architectural Study,'' \textit{GitHub Research Repository}, 2024.
\end{thebibliography}

\end{document}
